\chapter[Arbitraggio e top-level]{Arbitraggio e integrazione top-level}
\label{ch:top}

\section{\texttt{mem\_arbiter} --- arbitro a tre porte}
Un unico master di memoria byte-level (che alimenta la catena condivisa
\code{int8\_memory\_access} $\to$ \code{memory\_interface} $\to$ \code{psram\_controller})
è arbitrato tra tre richiedenti:

\begin{tabularx}{\textwidth}{C{1.3cm} L{3.4cm} Y}
\toprule
\rowh \thd{Porta} & \thd{Master} & \thd{Accessi} \\
\midrule
A & \code{spi\_engine} & \op{WRITE\_RAM} / \op{READ\_RAM}. \\
\rowa B & \code{neuron\_memory} & Letture X/W/bias durante un'esecuzione. \\
C & \code{layer\_sequencer} & Letture descrittori + scritture buffer tra layer. \\
\bottomrule
\end{tabularx}

Priorità fissa \textbf{B $>$ C $>$ A}: un'inferenza in corso è più critica della
contabilità del sequencer, che a sua volta è più critica di un accesso SPI manuale
appena arrivato. In funzionamento normale B e C sono comunque temporalmente disgiunti
(\code{neuron\_memory} richiede solo durante un'esecuzione, \code{layer\_sequencer} solo
nelle pause tra layer), quindi la priorità conta soprattutto per il caso limite di un
\op{WRITE\_RAM}/\op{READ\_RAM} manuale che arriva durante un'esecuzione multi-layer.

\begin{center}
\begin{tikzpicture}[font=\scriptsize,node distance=6mm]
 \node[fnblock,minimum width=30mm](a){Port A --- \code{spi\_engine}};
 \node[fnblock,below=4mm of a,minimum width=30mm](b){Port B --- \code{neuron\_memory}};
 \node[fnblock,below=4mm of b,minimum width=30mm](c){Port C --- \code{layer\_sequencer}};
 \node[fnblockD,right=16mm of b,minimum width=26mm,minimum height=16mm](arb){\code{mem\_arbiter}\\{\scriptsize B$>$C$>$A}};
 \node[fnblockT,right=14mm of arb,minimum width=26mm](m){catena memoria\\{\scriptsize condivisa}};
 \draw[fnarrow] (a)-|(arb.west|-a); \draw[fnarrow] (b)--(arb.west);
 \draw[fnarrow] (c)-|(arb.west|-c);
 \draw[fnbus] (arb)--(m);
\end{tikzpicture}
\end{center}

Concesso l'accesso, l'arbitro mantiene la proprietà fino all'impulso \code{m\_ready}
della singola transazione, poi rilascia: tutti e tre i master emettono \code{req} come
impulso pulito di un ciclo, quindi è sufficiente un design grant-and-forward senza code.

\section{\texttt{spi\_neuron\_top} --- integrazione completa}
Il top-level collega SPI (\code{spi\_slave}+\code{spi\_engine}), l'arbitro, il sequencer,
\code{neuron\_memory} e la catena PSRAM. Il reset di \code{neuron\_memory} è l'OR del
reset globale con l'impulso di soft-reset dell'opcode \op{RESET}, così l'host può
recuperare il motore via SPI senza reset fisico (la RAM resta intatta).

\begin{center}
\begin{tikzpicture}[font=\scriptsize,node distance=7mm]
 \node[fnblockA,minimum width=22mm](ss){\code{spi\_slave}};
 \node[fnblockA,right=8mm of ss,minimum width=22mm](se){\code{spi\_engine}};
 \node[fnblockT,below=8mm of se,minimum width=26mm](sq){\code{layer\_sequencer}};
 \node[fnblockD,right=10mm of se,minimum width=24mm](mux){MUX ctrl\\{\scriptsize su \code{seq\_busy}}};
 \node[fnblock,below=8mm of mux,minimum width=26mm](nm){\code{neuron\_memory}};
 \node[fnblockD,right=10mm of mux,minimum width=22mm](arb){\code{mem\_arbiter}};
 \node[fnblockA,right=8mm of arb,minimum width=26mm](mem){catena PSRAM};
 \draw[fnarrow] (ss)--(se);
 \draw[fnarrow] (se)--(mux);
 \draw[fnarrow] (sq)--(mux);
 \draw[fnarrow] (mux)--(nm);
 \draw[fnarrow] (se.south) to[bend right=10] (arb.north west);
 \draw[fnarrow] (nm)--(arb);
 \draw[fnarrow] (sq.east) to[bend right=20] (arb.south west);
 \draw[fnbus] (arb)--(mem);
\end{tikzpicture}
\end{center}

Il multiplexer commuta le linee di controllo di \code{neuron\_memory} tra il sequencer
(mentre \code{seq\_busy} è alto) e il percorso diretto di \code{spi\_engine} (modalità
single-layer legacy), restituendo il motore al percorso diretto a fine sequenza.

\begin{fnnote}[Verifica end-to-end]
\code{spi\_neuron\_top} è verificato in simulazione con PSRAM reale
(\code{psram\_model.v}, nessun mock): RESET/READ\_CONFIG/WRITE\_RAM/READ\_RAM/SET\_BASE/
START/STATUS/READ\_OUTPUT e \op{RUN\_NETWORK} sono esercitati puramente su SPI simulato
(cap.~\ref{ch:impl}).
\end{fnnote}
